\section{TEE + Untrusted-GPU Split Inference: What Holds}
\label{sec:tee}

% LOCKED campaign: 1 modern model x 1 TEE platform (e.g. TDX + H100 CC).
% Deliverables: (a) latency/throughput/overhead breakdown + bottleneck analysis;
% (b) the SAME attack suite as AloePri, run against the split, showing it resists.
% Status: IN PROGRESS (co-developing with the paper).

\subsection{Architecture}
% Boundary layers (embedding lookup, lm_head / pooling) inside the TEE; transformer
% blocks on the untrusted GPU over obfuscated activations; inverse applied before
% non-linearities. State the exact split for the chosen model.
\needswork{describe the split for the chosen model + platform}
\dnote{open decision: exact model + TEE platform — stage-1 still-open item \#3}

\subsection{Implementation}
\needswork{implementation notes: what runs where, attestation/RATLS setup}

\subsection{Performance: Overhead and Bottlenecks}
% The systems result. Latency, throughput, overhead vs plaintext baseline; where
% the time goes (attestation, data movement, inverse ops, GPU compute).
\begin{table}[t]
  \centering
  \caption{Split-inference overhead vs.\ plaintext baseline.}
  \label{tab:tee-overhead}
  \begin{tabular}{@{}lrrr@{}}
    \toprule
    Configuration & Latency & Throughput & Overhead \\
    \midrule
    Plaintext baseline & \needswork{} & \needswork{} & --- \\
    TEE+GPU split      & \needswork{} & \needswork{} & \needswork{} \\
    \bottomrule
  \end{tabular}
\end{table}
\needswork{fill overhead table + bottleneck breakdown figure}

\subsection{Attack Resistance}
% Run inversion / GRAM / ISA AttnScore against the obfuscated activations the GPU
% sees. Show attack success is low — this is the "holds" evidence and must use the
% SAME harness/metrics as Sec. 5 for a fair comparison (see Sec. 6).
\needswork{attack-success numbers on the split; same metrics as AloePri}
